% TOP_PROBLEM: {"id": 32, "title": "The Kakeya Conjecture", "category_id": 6, "category": {"id": 6, "name": "geometry", "display_name": "Geometry", "description": "Euclidean and non-Euclidean geometry, geometric structures.", "slug": "geometry", "order_index": 6, "created_at": "2026-07-31T15:26:25.671Z"}, "status": "open"}
% FIRST_POSTED: 2026-10-06
% ATTEMPT_STATUS: unresolved
% OPENAI_CLAIM: {"schema_version":1,"source_repo":"https://github.com/openai/math","source_commit":"adc7f1241b42e322a6451854ab7e4b4c146bf78a","release_date":"2026-10-06","match":"partial","resolution":"partial","independently_reviewed":false,"second_pass":"n/a","families":[{"family":"074","title":"Kakeya in three and four dimensions","subject":"Real and complex analysis","manuscripts":[{"title":"Every four-dimensional Kakeya set has full Hausdorff dimension","pdf_path":"preprints/Every-four-dimensional-Kakeya-set-has-full-Hausdorff-dimension-September-24-2026/paper.pdf","readme_path":"preprints/Every-four-dimensional-Kakeya-set-has-full-Hausdorff-dimension-September-24-2026/README.md","date":"2026-09-24","theorem_ref":"Theorem/corollary (label thm:main), preprints/Every-four-dimensional-Kakeya-set-has-full-Hausdorff-dimension-September-24-2026/build/sections/01-introduction.tex"},{"title":"The Kakeya maximal conjecture in three dimensions","pdf_path":"preprints/The-Kakeya-maximal-conjecture-in-three-dimensions-September-23-2026/paper.pdf","readme_path":"preprints/The-Kakeya-maximal-conjecture-in-three-dimensions-September-23-2026/README.md","date":"2026-09-23","theorem_ref":"Theorem/corollary (label thm:main), preprints/The-Kakeya-maximal-conjecture-in-three-dimensions-September-23-2026/build/sections/introduction.tex"}],"lean":null,"reasoning_trace":null}]}
\documentclass[11pt]{article}
\usepackage[margin=25mm]{geometry}
\usepackage{fontspec}
\setmainfont{Latin Modern Roman}
\usepackage{amsmath,amssymb,mathtools}
\usepackage{unicode-math}
\setmathfont{Latin Modern Math}
\usepackage{enumitem}
\usepackage{xurl,hyperref,bookmark}
\hypersetup{colorlinks=true,urlcolor=blue}
\setcounter{secnumdepth}{0}
\setlength{\parindent}{0pt}
\setlength{\parskip}{5pt}
\setlength{\emergencystretch}{3em}
\begin{document}
\section*{The Kakeya Conjecture}

\subsection{OpenAI's claim}
\textit{Four-dimensional Kakeya sets}: The manuscript claims that a set in $\mathbb R^4$ containing a segment of length one in every direction has Hausdorff dimension four; it permits arbitrary sets and arbitrary choices of witnessing segments.

\textit{Three-dimensional Kakeya maxima}: The claimed tube maximal estimate in three dimensions has $L^3$ input and output and loses only an arbitrarily small power of the tube radius; in particular, the density exponent in its corresponding shaded-tube estimate is three.

\subsection{Outline of the argument}
\textit{Four-dimensional Kakeya sets}: A hypothetical dimension loss is converted into a weighted incidence model for lines. Nested boxes and quadratic fits describe persistent concentration; each fit must retain enough incidence mass to contribute. Scalar comparison and planar rigidity then contradict the simultaneous density and time-length requirements forced on substantial terminal chart systems.

\textit{Three-dimensional Kakeya maxima}: The proof keeps the fraction of each tube that is actually observed throughout a multiscale concentration model. Near-extremal indexed configurations are made stationary at critical exponents, while their position and direction information are retained jointly. The geometric argument excludes a positive critical loss, and the final reduction passes from tube shadings to general functions without worsening the cubic density dependence.

\subsection{Scope relative to this problem}
\textit{Four-dimensional Kakeya sets}: The paper resolves a full-direction Kakeya statement only in its stated dimension (three for the maximal estimate, four for the set theorem). The question quantifies over every dimension, or over arbitrary direction sets with the dimension(D)+1 bound, leaving those other cases open.

\textit{Three-dimensional Kakeya maxima}: The paper resolves a full-direction Kakeya statement only in its stated dimension (three for the maximal estimate, four for the set theorem). The question quantifies over every dimension, or over arbitrary direction sets with the dimension(D)+1 bound, leaving those other cases open.

This page records partial support. Complete resolution of the catalogue question is not established by this record.

\subsection{Formal verification}
\textit{Four-dimensional Kakeya sets}: The pinned catalogue lists no Lean formalization for this result.

\textit{Three-dimensional Kakeya maxima}: The pinned catalogue lists no Lean formalization for this result.

This site has not run the Lean development. These links describe OpenAI’s selected formal statements and should be read with their stated scope.

\subsection{Status caveat}
OpenAI reports that an internal model produced the manuscripts. Their arguments have not been independently verified as proofs by this site. Any portion outside the documented formal scope may contain mistakes, and the catalogue label reports the announced mathematical claim.

\subsection{Sources}
\begin{itemize}
\item OpenAI, \emph{Every four-dimensional Kakeya set has full Hausdorff dimension}, 2026-09-24; citation key \texttt{OAI:Every-four-dimensional-Kakeya-set-has-full-Hausdorff-dimension-September-24-2026}. \href{https://github.com/openai/math/blob/adc7f1241b42e322a6451854ab7e4b4c146bf78a/preprints/Every-four-dimensional-Kakeya-set-has-full-Hausdorff-dimension-September-24-2026/paper.pdf}{PDF (GitHub)}; \href{https://github.com/openai/math/raw/adc7f1241b42e322a6451854ab7e4b4c146bf78a/preprints/Every-four-dimensional-Kakeya-set-has-full-Hausdorff-dimension-September-24-2026/paper.pdf}{PDF (direct)}; \href{https://github.com/openai/math/blob/adc7f1241b42e322a6451854ab7e4b4c146bf78a/preprints/Every-four-dimensional-Kakeya-set-has-full-Hausdorff-dimension-September-24-2026/README.md}{Manuscript page and citation}; \href{https://github.com/openai/math/blob/main/preprints/Every-four-dimensional-Kakeya-set-has-full-Hausdorff-dimension-September-24-2026/paper.pdf}{Latest version}.
\item OpenAI, \emph{The Kakeya maximal conjecture in three dimensions}, 2026-09-23; citation key \texttt{OAI:The-Kakeya-maximal-conjecture-in-three-dimensions-September-23-2026}. \href{https://github.com/openai/math/blob/adc7f1241b42e322a6451854ab7e4b4c146bf78a/preprints/The-Kakeya-maximal-conjecture-in-three-dimensions-September-23-2026/paper.pdf}{PDF (GitHub)}; \href{https://github.com/openai/math/raw/adc7f1241b42e322a6451854ab7e4b4c146bf78a/preprints/The-Kakeya-maximal-conjecture-in-three-dimensions-September-23-2026/paper.pdf}{PDF (direct)}; \href{https://github.com/openai/math/blob/adc7f1241b42e322a6451854ab7e4b4c146bf78a/preprints/The-Kakeya-maximal-conjecture-in-three-dimensions-September-23-2026/README.md}{Manuscript page and citation}; \href{https://github.com/openai/math/blob/main/preprints/The-Kakeya-maximal-conjecture-in-three-dimensions-September-23-2026/paper.pdf}{Latest version}.
\item \href{https://github.com/openai/math/blob/adc7f1241b42e322a6451854ab7e4b4c146bf78a/CONTENTS.md}{OpenAI catalogue entry}.
\end{itemize}
\end{document}
